\clearpage\fancyhead[L]{\small PALLAS.}\hypertarget{native-325-lower}{}\label{native-325-lower}
\begin{center}\textbf{PALLAS.}\par\vspace{4pt}\hrule\vspace{4pt}\fontsize{10}{12}\selectfont Monatliche Correspondenz. Band XXII. S. 591\,..\,595. 1810 December.\par\vspace{4pt}\hrule\vspace{4pt}\end{center}

\par\indent „Ich eile, schrieb Herr Professor \textsc{Gauss}, Ihnen theuerster Freund, eine wichtige Nachricht mitzutheilen. Ich schrieb Ihnen schon neulich, dass ich mit Berechnungen der Störungen beschäftigt sei, die Pallas in den Jahren 1803\,..\,1811 vom Jupiter erlitten hatte. Wie wenig rein elliptische Elemente alle bisher beobachteten sechs Oppositionen darzustellen vermögen, habe ich in der Abhandlung gezeigt, deren ich schon einigemal gegen Sie erwähnt habe. Es war also, nachdem ich vor einigen Tagen die Arbeit über jene Störungen vollendet hatte, eine äusserst angenehme Satisfaction für mich, zu finden, dass nun mit Berücksichtigung jener Störungen alle sechs Oppositionen über meine Erwartung genau sich vereinigen liessen. Hier das Tableau der Abweichungen:

\begin{center}\fontsize{10}{12}\selectfont\setlength{\tabcolsep}{4pt}\renewcommand{\arraystretch}{1.15}\begin{tabular}{r|r|r}
 & \shortstack{Mittlere\\Länge} & \shortstack{Heliocentr.\\Breite}\\\hline
1803 & $+1''3$ & $-1''0$\\
1804 & $-3.8$ & $+4.1$\\
1805 & $+3.9$ & $+6.0$\\
1807 & $-3.3$ & $+3.9$\\
1808 & $+3.2$ & $-17.0$\\
1809 & $-1.4$ & $-3.9$\\
\end{tabular}\end{center}
\noindent (Die Breite 1808 war bekanntlich sehr schlecht bestimmt.)

\SourceLeaf{326}\par\indent „Da 1809 bei der heliocentrischen Conjunction mit dem Jupiter die Elemente sich \emph{sehr stark} geändert hatten, so war ich neugierig zu sehen, wie viel meine Ephemeride im October-Hefte der \emph{Monatl. Corresp.} wohl fehlen könne, und berechnete nach meinen neuen Resultaten mit Rücksicht auf die Störungen einen Ort. Hier fand ich nun für 1811 den 9. Januar Mittag in Göttingen

{\fontsize{10}{12}\selectfont\[\begin{array}{llr|l}\text{AE }\PallasSymbol & 151^\circ\;37' & \text{also }1^\circ\;40'\text{ weniger} & \raisebox{-0.6\baselineskip}{als in der Ephemeride.}\\\text{Decl.} & 22^\circ\;10'\text{ südl.} & \text{—}\quad 11'\text{ südl.} & \end{array}\]}

\noindent Ich habe jetzt eben noch einen zweiten Ort berechnet, 1810 den 26. October Mittag in Göttingen

{\fontsize{10}{12}\selectfont\[\begin{array}{llr|l}\text{AE }\PallasSymbol & 134^\circ\;55' & \text{also }45'\text{ weniger} & \raisebox{-0.6\baselineskip}{als die Ephemeride}\\\text{Decl.} & 14^\circ\;28'\text{ südl.} & \text{—}\quad 6'\text{ südl.} & \end{array}\]}

\Needspace{15\baselineskip}\noindent Wünschen Sie selbst einzelne Oerter zu berechnen, so können Sie sich dazu folgender Elemente bedienen:

\begin{center}\fontsize{10}{12}\selectfont\setlength{\tabcolsep}{4pt}\renewcommand{\arraystretch}{1.15}\begin{tabular}{llr}
Epoche 1811 Göttingen & .\;.\;.\;.\;.\;. & $126^\circ\;32'\;52''$\\
tägliche mittlere trop. Bewegung & .\;. & $769''012$\\
Logarithmus der halben Axe & .\;.\;.\;. & $0.44277$\\
Excentricität & .\;.\;.\;.\;.\;.\;.\;.\;.\;. & $0.24162\ (=\operatorname{Sin}13^\circ\;18'\;56'')$\\
Perihelium für ob. Epoche & .\;.\;.\;.\;. & $120^\circ\;55'\;5''$\\
Knoten & .\;.\;.\;.\;.\;.\;.\;.\;.\;.\;.\;. & $172\;32\;30$\\
Neigung & .\;.\;.\;.\;.\;.\;.\;.\;.\;.\;. & $34\;35\;55$\\
\end{tabular}\end{center}

\noindent Noch einen Ort kann ich vor Abschickung hinzufügen 1810 den 15. Dec. Mittag in Göttingen.

{\fontsize{10}{12}\selectfont\[\begin{array}{llr|l}\text{AE }\PallasSymbol & 149^\circ\;37' & \text{also }1^\circ\;15'\text{ weniger} & \raisebox{-0.6\baselineskip}{als die Ephemeride}\\\text{Decl.} & 21''\;37'\text{ südl.} & \text{—}\quad 12'\text{ südlicher} & \end{array}\]}

\par\indent „Ich bitte Sie nun, die erstere heitere Nacht zur Aufsuchung der Pallas anzuwenden. Hier ist es bisher immer unmöglich gewesen. Mit welcher Ungeduld ich die Bestätigung dieses merkwürdigen Resultates erwarte, das bisher bei den neuen Planeten noch nicht vorgekommen ist, können Sie leicht denken.

\par\indent „Die Verbesserung der Ephemeride im Octoberheft durch Interpolation ist folgende:
\Needspace{12\baselineskip}
\begin{center}\fontsize{10}{12}\selectfont\setlength{\tabcolsep}{4pt}\renewcommand{\arraystretch}{1.15}\begin{tabular}{ll|r|l}
 &  & AE. & Decl.\\\hline
1810 den 26. October & Mittag & $-45'$ & $6'$ südl.\\
--- den 7. November & Mitternacht & 49 & 8\\
--- den 20. November & Mittag & 56 & 10\\
--- den 2. December & Mitternacht & 65 & 11\\
--- den 15. December & Mittag & 75 & 12\\
--- den 27. December & Mitternacht & 87 & 12\\
1811 den 9. Januar & Mittag & 101 & 11\\
\end{tabular}\end{center}

\par\indent So merkwürdige Resultate und der Wunsch unseres verehrtesten Freundes, waren Aufforderung genug, um mit möglichstem Fleisse die Aufsuchung der Pallas zu versuchen. Leider vereitelte der seit sechs Wochen höchst ungünstige Himmel fast alle unsere Bemühungen. Am 18. December wurde es etwas helle, und im parallactischen Instrument sahen wir Morgens zwei Uhr drei Sterne im Felde des Fernrohrs, von denen einer höchst wahrscheinlich die Pallas war; allein eine bestimmte Beobachtung wurde unmöglich, da sich der Himmel bald wieder umzog. Auch haben wir von jenen drei Sternen nirgends weder im \emph{Piazzi} noch in der \emph{Histoire céleste} eine Bestimmung auffinden können. Am 19. De\SourceLeaf{327}cember gelang uns endlich eine Meridian-Beobachtung. Genau an der Stelle, wo die verbesserte Ephemeride den Pallas-Ort gab, beobachteten wir einen Stern 7...8. Grösse, welcher höchst wahrscheinlich die Pallas war.

\begin{center}\fontsize{10}{12}\selectfont\setlength{\tabcolsep}{4pt}\renewcommand{\arraystretch}{1.15}\begin{tabular}{c|c|c|c}
1809 & \shortstack{mittlere\\Seeberger Zeit} & AR. app. $\PallasSymbol$ & \shortstack{Declination\\austr.}\\\hline
December 19 & $16^{\text{u}}\;9^{\text{m}}\;7^{\text{s}}3$ & $150^\circ\;20'\;13''5$ & $22^\circ\;58''7$\\
\end{tabular}\end{center}

\par\indent Die Declination konnte freilich nur am Passage-Instrument beobachtet werden, und ist auf die Minute ungewiss. Die AR. dagegen ist genau und diese stimmt mit der verbesserten Bestimmung ganz vortrefflich überein. Im Fernrohr des Quadranten war Pallas durchaus nicht sichtbar. Ein kleiner in der Nähe der Pallas befindliche Stern, dessen Ort wir fanden
\[\text{AR. }149''\;28'\;36''.\quad\text{Decl. austr. }21^\circ\;59'\]
\noindent hatte unverrückt seine Lage behalten. Wir standen an, diese isolirte Beobachtung Herrn Prof. \textsc{Gauss} mitzutheilen, in der Hoffnung, bald durch eine zweite es constatiren zu können, dass es wirklich Pallas war. Allein leider gestattete der Himmel uns diese Freude bis jetzt nicht, und wir theilen daher ersterm unsern vermutheten Pallas-Ort mit. \textsc{Gauss} der unserer Vermuthung beitritt, schrieb uns darauf folgendes:

\par\indent „Sehr freue ich mich, dass meine Rechnung über die Pallas-Störungen sich so schön belohnt. Denn obgleich bisher das Wetter es mir unmöglich gemacht hat, den Planeten aufzusuchen, so zweifle ich doch nicht mehr, dass der von Ihnen beobachtete Stern die Pallas gewesen ist. \textsc{Harding} hat auf seiner Karte an dem Platz keinen Stern, aber den andern, welchen Sie am 18. Dec. mit sahen, hat er, obwohl er in der \emph{Histoire céleste} nicht vorkömmt, angemerkt.“
